\begin{abstract}
This study proposes a temporal modeling framework with a counterfactual policy-simulation layer to examine student dropout in higher education using time-stamped learning-management-system (LMS) engagement data and administrative withdrawal records. Dropout is operationalized as an observed time-to-event outcome at the course-run enrollment level, and weekly risk is modeled in discrete time via a penalized, class-balanced logistic regression over weekly person--period rows. As a robustness check, we also fit a Cox proportional hazards model with time-varying covariates.

Under a late-event temporal holdout, the discrete-time model attains row-level AUCs of 0.8359 on training and 0.8396 on test for discriminating weekly hazard. Complementary ablation analyses indicate that performance and survival diagnostics are sensitive to the feature set composition, highlighting the role of temporal engagement signals in risk modeling. For decision-oriented evaluation, a structural hazard-scaling policy is examined via counterfactual simulation, producing survival contrasts $\Delta S(T)$ under explicit, auditable scenario parameters. Beyond aggregate evaluation, the framework is applied in a subgroup-sensitive analysis to compare how the same scenario changes survival gaps between groups, with uncertainty quantified via bootstrap. These results are not presented as causally identified effects, but as evidence that the mathematical framework supports structural policy evaluation and consistent scenario comparison when intervention effects are not identifiable from the available observational data.
\end{abstract}